\section{Technical Overview}

Our goal is to share a sparse map without revealing its indices.
The parties first supply shares of at most $t$ addresses.
Later, they supply new payloads for those same addresses.
Each expansion returns shares of a length-$N$ map; duplicate addresses
contribute their payload sum.
A dummy address contributes nothing.
The ideal functionality lets the corrupt party choose its output share;
the honest party receives the complement to the specified map.
The main cost is local expansion: a sum of $t$ independent DPFs traverses
the full domain $t$ times.

\paragraph{Sparse DPFs.}
A single-point DPF uses a tree of pseudorandom seeds and public correction
words to keep one secret path active. For a public subset $S\subseteq[0,N)$,
our sparse DPF contracts paths with only one live child.
It materializes only the leaves in $S$ and their branching nodes.
Setup expands this compressed tree and saves the leaf seeds.
Each later expansion converts those seeds using a round-indexed generator
and applies a payload correction. It processes $|S|$ leaves rather than
traversing the full ambient domain.

\paragraph{The cuckoo layout.}
A cuckoo-based DMPF gives each address one candidate bin in each of $w$
partitions \cite{fssBatchCodes,cuckooDmpf}.
Each bin owns the public sparse set of addresses that can map to it.
If the secret input rows receive distinct admissible bins, one sparse DPF
per bin suffices. Summing the bin outputs reconstructs the multi-point map.
Each address appears once per partition, so expansion visits $wN$ leaves.
The difficulty is finding the secret placement efficiently in MPC.

\subsection{Waterfall Cuckoo DMPF}

Waterfall first merges duplicate addresses and pads the list with dummies.
It then samples public $t$-wise independent hashes before routing.
Distinct real addresses have independent uniform candidates; dummy rows
receive fresh uniform candidates.
Consequently, the entire candidate matrix has the same distribution for
every input support.

An inexpensive scan fills the partitions in order.
A row that collides remains unplaced and tries the next partition.
A small public number of complete augmentations then repairs the remaining
matching, followed by stash compaction.

Setup hides the resulting placement with two privately controlled permutation
passes and opens only the permuted placement.
On successful routing, that opening is a uniform labeled injection.
It supports an oblivious scatter into the sparse-DPF columns.
The simulation proof keeps the hashes unconditioned and charges routing
failure once per setup state.
It uses ideal MPC, permutation, and reusable-DPF functionalities;
concrete realizations contribute their separate security losses.

At $t=16$, the default profile uses four 16-bin partitions and three repairs,
with no stash. It has $4N$ local leaf expansion.
The lower-expansion profile uses partition sizes $(128,128,64)$ and two
repairs, also without a stash. It saves one domain traversal per expansion
but has larger setup and state.

\subsection{Bucketing DMPF}

For independent uniform input addresses, public intervals can divide the
inputs into smaller padded groups. One DMPF handles each group.
The tradeoff depends on the grouping circuit, the underlying DMPF, and an
explicit overflow bound. Expansion must still materialize $N$ outputs.

\subsection{Big-State DMPF}

The big-state approach \cite{dmpf} trades stored state for cheaper subsequent
evaluation. We discuss a black-box key-generation path and a generic-MPC
path for this setup--state tradeoff.

\subsection{Integration into Ring-LPN and Stationary-LPN PCGs}

A Ring-LPN PCG must share products of sparse polynomials.
With $P$ polynomials per party and $t$ nonzero coefficients per polynomial,
each cross-party product has $t^2$ terms before cancellation.
For regular noise, one coefficient lies in each block of length
$L=N_{\mathsf R}/t$.
The products reorganize into $P^2t$ lists, each with $t$ terms and a
local domain of length $2L$.
Replacing a sum of $t$ DPFs on each list by a constant-expansion DMPF removes
a factor of $t$ from the dominant leaf work.

Stationary supports allow the address-dependent setup to be reused.
Every expansion samples fresh coefficients and a fresh public mask vector.
It converts cached leaves, communicates payload corrections, and performs
ring arithmetic. Thus repeated expansion is not communication-free.
Security across a polynomial number of expansions requires a joint
stationary Ring-LPN assumption, separate from the routing argument.

\subsection{Reverse Cuckoo}

For these random Ring-LPN supports, we also consider a placement-first
construction. Reverse Cuckoo chooses a hidden placement, then programs
public hashes compatible with it. This gives the $2N$ specialized pipeline
evaluated in the paper.
Its descriptors leak information about the support: in the ideal hash model,
a hash tuple is weighted by its number of valid placements.
Hidden permutations do not remove this multiplicity bias.
We therefore use Reverse Cuckoo under an explicit leakage-robust decisional
Ring-LPN assumption. The degree-two evaluator's correctness analysis is
separate from that assumption. \sectionref{sec:revCuckoo} derives the leakage
distribution and examines attacks that use it.
